% Source: https://artofproblemsolving.com/community/c3h2915821_2022_mbmt_team_round__montgomery_blair_math_tournament
% Stable identity: aops-topic-2915821-q12
In a regular pentagon, $\angle CDE=108^\circ $. Since $\angle CDF=60^\circ $, $\angle FDE=48^\circ $. Also $ DE=DC=DF $, so isosceles triangle $ DEF $ has base angle $\angle DEF=66^\circ $. The diagonal $ EB $ makes $\angle BED=72^\circ $; for example, the intercepted arc $ BD $ through $ C $ is $ 144^\circ $, giving half that inscribed angle. Therefore $\angle BEF=72^\circ-66^\circ=\boxed{6^\circ}$.
\begin{asy}
usepackage("amsmath");
size(7cm);
pair C=(-0.5,0),D=(0.5,0),PtE=D+dir(72),A=PtE+dir(144),B=A+dir(216),F=(0,sqrt(3)/2);draw(A--B--C--D--PtE--cycle);draw(C--F--D^^PtE--B^^PtE--F,blue);label("$A$",A,N);label("$B$",B,W);label("$C$",C,SW);label("$D$",D,SE);label("$E$",PtE,E);label("$F$",F,S);
\end{asy}
